\DocumentMetadata{
  pdfversion=2.0,
  pdfstandard=ua-2,
  lang=en-US,
  tagging=on,
  tagging-setup={math/setup=mathml-SE,tabsorder=structure}
}
\documentclass[11pt,letterpaper]{article}
\usepackage[margin=0.68in,includefoot]{geometry}
\usepackage{newcomputermodern}
\usepackage{amsmath,amscd,mathtools}
\usepackage{tikz,adjustbox}
\providecommand{\square}{\mdlgwhtsquare}
\usepackage[hidelinks]{hyperref}
\hypersetup{
  pdftitle={Analysis PhD Exam, May 1991},
  pdfauthor={University of Florida Department of Mathematics},
  pdfsubject={Graduate mathematics examination},
  pdfdisplaydoctitle=true,
  bookmarksopen=true
}
\setcounter{secnumdepth}{0}
\setlength{\parindent}{0pt}
\setlength{\parskip}{0.28em}
\setlength{\emergencystretch}{3em}
\setlength{\leftmargini}{2.15em}
\setlength{\leftmarginii}{2.65em}
\setlength{\leftmarginiii}{3em}
\pagestyle{empty}
\raggedbottom
\allowdisplaybreaks
\NewTaggingSocketPlug{block/list/label}{examproblem}{%
  \tagstructbegin{tag=Lbl}\tagstructend%
  \tagstructbegin{tag=\LBody}%
  \tagstructbegin{tag=\ProblemHeadingTag,title-o={\ProblemHeadingText}}%
  \tagmcbegin{tag=\ProblemHeadingTag,actualtext={\ProblemHeadingText}}%
  #2%
  \tagmcend%
  \tagstructend%
}
\newenvironment{examproblems}[2]{%
  \def\ProblemHeadingTag{#2}%
  \AssignTaggingSocketPlug{block/list/label}{examproblem}%
  \begin{enumerate}%
  \setcounter{enumi}{#1}%
  \renewcommand{\labelenumi}{\textbf{\theenumi.}}%
  \setlength{\topsep}{0.35em}%
  \setlength{\partopsep}{0pt}%
  \setlength{\itemsep}{0.48em}%
  \setlength{\parsep}{0.18em}%
}{\end{enumerate}}
\newenvironment{examsubparts}{%
  \AssignTaggingSocketPlug{block/list/label}{default}%
  \begin{enumerate}%
  \setlength{\topsep}{0.18em}%
  \setlength{\partopsep}{0pt}%
  \setlength{\itemsep}{0.16em}%
  \setlength{\parsep}{0.1em}%
}{\end{enumerate}}
\newenvironment{examinstructions}{%
  \par\begingroup\itshape
}{%
  \par\endgroup\vspace{0.18em}
}
\makeatletter
\renewcommand\subsection{\@startsection{subsection}{2}{0pt}%
  {-1.25ex plus -.25ex minus -.1ex}%
  {0.55ex plus .1ex}%
  {\normalfont\large\bfseries}}
\renewcommand{\@maketitle}{%
  \newpage\null\vskip -1em%
  {\footnotesize\bfseries
    \href{https://www.ufl.edu/}{University of Florida}, \href{https://math.ufl.edu/}{Department of Mathematics}\par}%
  \vskip -0.5em%
  \tagstructbegin{tag=H1,title={Analysis PhD Exam, May 1991}}%
  \tagpdfparaOff%
  \tagmcbegin{tag=H1}%
    {\LARGE\bfseries\href{https://gma.math.ufl.edu/exams/analysis-phd/}{Analysis PhD Exam}, May 1991\par}%
  \tagmcend%
  \tagpdfparaOn%
  \tagstructend%
  \vskip 0.25em%
  \tagmcbegin{artifact}%
    \hrule height 1pt%
  \tagmcend%
  \vskip 1em%
}
\makeatother
\title{Analysis PhD Exam, May 1991}
\author{}
\date{}
\begin{document}
\maketitle
\thispagestyle{empty}
\begin{examinstructions}
Be sure to give complete proofs and justify your steps. Do not leave gaps.
\end{examinstructions}
\begin{examproblems}{0}{H2}
\def\ProblemHeadingText{Problem 1.}
\item
State:
\begin{examsubparts}
\item[\textnormal{(a)}]
The Lebesgue Dominated Convergence Theorem
\item[\textnormal{(b)}]
Fatou’s Lemma
\item[\textnormal{(c)}]
The Radon–Nikodym Theorem
\item[\textnormal{(d)}]
\((L^1(S,\Sigma,\mu))^* = L^\infty(S,\Sigma,\mu)\)
\end{examsubparts}
\def\ProblemHeadingText{Problem 2.}
\item
State and prove the Vitali Convergence Theorem for integrals (convergence in measure form).
\def\ProblemHeadingText{Problem 3.}
\item
Show that a subset~\(A\) of the reals is a set of Lebesgue measure zero if and only if there exists a sequence of intervals \(I_m\) such that
\begin{examsubparts}
\item[\textnormal{(a)}]
\(\displaystyle \sum_m \operatorname{length}(I_m)<\infty.\)
\item[\textnormal{(b)}]
\(\displaystyle A\subset \bigcup_{m=k}^{\infty} I_m\) for every~\(k\).
\end{examsubparts}
\def\ProblemHeadingText{Problem 4.}
\item
Let~\(f\) be Lebesgue integrable on \(\mathbb R\). Suppose that \(\displaystyle \int_0^x f\,dm=0\) for every \(x\in\mathbb R\). What can you conclude about~\(f\)?
\def\ProblemHeadingText{Problem 5.}
\item
State and prove Fubini’s Theorem.
\def\ProblemHeadingText{Problem 6.}
\item
Let \((S_m,\Sigma_m,P_m)_{m=1}^{\infty}\) be a sequence of probability spaces. Let \(E_m\in\Sigma_m\), \(m=1,2,\ldots\).
\begin{examsubparts}
\item[\textnormal{(a)}]
Show that
\[
E=\mathop{\times}_{m=1}^{\infty}E_m\in\bigotimes_{m=1}^{\infty}\Sigma_m.
\]
\item[\textnormal{(b)}]
Prove that
\[
\left(\bigotimes_{m=1}^{\infty}P_m\right)(E)=\prod_{m=1}^{\infty}P_m(E_m).
\]
\end{examsubparts}
\end{examproblems}
\end{document}
